\section{Recomendaciones y Selección Recomendada}

Con base en el análisis comparativo de las Secciones 7 y 8, se recomienda la siguiente selección de equipos y componentes:

\begin{enumerate}
\item Ventilador axial mural (placa mural) Ø560 mm en PRFV o ejecución epóxica anticorrosiva con transmisión directa, uniforme con el montaje típico instalado en la planta: banco de filtración alojado en la cubierta intemperie, estructura de unión pernada al muro y malla de protección interior. El motor va en el cubo del impulsor dentro de la corriente de aire, en ejecución encapsulada severe duty; la alternativa de transmisión por bandas (motor fuera de la corriente) queda documentada como opción. Primera opción de referencia: Sodeca HQD/HGT mural anticorrosivo (filial local); alternativas: Greenheck mural, Aerovent/Twin City mural FRP y New York Blower FRP mural. Motor provisional 0.75 HP TEFC encapsulado severe duty anticorrosivo, 440 V, 3$\phi$, 60 Hz.
\item Filtración en dos etapas: prefiltro MERV 8 y filtro final V-bank MERV 13-14 de marco plástico (Camfil Durafil ES2 MERV 14 como referencia de diseño; AAF VariCel VXL y Freudenberg MaxiPleat MX 85 como alternativas), con portafiltros inox 304/316. Se recomienda congelar la especificación en términos MERV/ePM1 y dimensiones 24$\times$24 in para alternar proveedores de repuestos.
\item Tres rejillas eggcrate inox 316L de 353 mm $\times$ 336 mm con malla anti-insectos inox 316 18$\times$18 desmontable, fabricación local a medida (ProAire/Laminaire).
\item Ferretería inox 316 (A4), sellante silicona RTV neutra y aislamiento dieléctrico frente a estructuras galvanizadas en todos los accesorios de montaje (estructura de unión, cubierta intemperie y malla de protección interior).
\item Confirmar con los proveedores, antes de la compra, el tamaño y RPM exactos del ventilador mediante el software del fabricante para el punto de catálogo 3 840 m$^3$/h a 225 Pa, y verificar que la potencia del motor cubra el escenario de filtro cargado con el margen de servicio previsto.
\item Verificar el caudal real en el ensayo de balanceo mediante anemometría en las tres rejillas de descarga, comparándolo con el caudal de diseño de 3 840 m$^3$/h.
\end{enumerate}
